%% kapa: exam=1 round=27 year=2016 period=2 form=A subject=appraisal_law subject_ko=감정평가관계법규 seq=27 number=27
%% source: Q-Net 감정평가사 기출문제 (cst00309) artlSeq=5169402 file=extracted/1008/5169402/제27회 감정평가사 1차 문제/제27회 감정평가사 2교시문제(A형).pdf
%% conversion: kordoc 4.24.51 -> tools/gen1.py (auto); review=proofread-sonnet; flags=none
\begin{question}{27}
건축법령상「행정대집행법」적용의 특례 사유로 규정되지 않은 것은?
\begin{choices}
  \item 어린이를 보호하기 위하여 필요하다고 인정되는 경우
  \item 재해가 발생할 위험이 절박한 경우
  \item 건축물의 구조 안전상 심각한 문제가 있어 붕괴 등 손괴의 위험이 예상되는 경우
  \item 허가권자의 공사중지명령을 받고도 불응하여 공사를 강행하는 경우
  \item 도로통행에 현저하게 지장을 주는 불법건축물인 경우
\end{choices}
\end{question}
\answer{1}
